%----------------------------------------------------------------------------
% Course Syllabus for Chemistry 352 (c) 2018 by Dale J. Brugh

% Course Syllabus for Chemistry 352 is released under a
% Creative Commons Attribution-ShareAlike 4.0 International License.

% See http://creativecommons.org/licenses/by-sa/4.0/ for
% a description of your rights under this license. 
%----------------------------------------------------------------------------
\documentclass[letterpaper,oneside,onecolumn,11pt,article]{memoir}
%
% --- LOAD PACKAGES ---
%
\usepackage[T1]{fontenc}            % use T1 font encoding
\usepackage{textcomp}
\usepackage{courier}                % set courier as typewriter font
\usepackage{times}                  % set times as text font
\usepackage[scaled=0.92]{helvet}    % set Helvetica as the sans-serif font
\usepackage{mtpro2}
\usepackage{setspace}
\usepackage{amsmath}
\usepackage{graphicx,color}
\usepackage{wallpaper}
\usepackage{textcomp}
\usepackage{relsize,fancyvrb}
\usepackage{verbatim}
\usepackage{caption}
\usepackage{paralist}
\usepackage{boxedminipage}
\usepackage{enumerate}
\usepackage[bookmarks=true]{hyperref}
%
% --- HYPER SETUP ---
%
\hypersetup{
    unicode=false,          % non-Latin characters in Acrobat’s bookmarks
    pdftoolbar=true,        % show Acrobat’s toolbar?
    pdfmenubar=true,        % show Acrobat’s menu?
    pdffitwindow=true,      % page fit to window when opened
    pdftitle={Syllabus: Chemistry 352 / Spring 2018}, 
    pdfauthor={Dale J. Brugh},     % author
    pdfsubject={Physical Chemistry},   % subject of the document
    pdfnewwindow=true,      % links in new window
    pdfkeywords={classes, ch352s18}, % list of keywords
    colorlinks=true,       % false: boxed links; true: colored links
    linkcolor=black,          % color of internal links
    citecolor=green,        % color of links to bibliography
    filecolor=magenta,      % color of file links
    urlcolor=black           % color of external links
}
\definecolor{nicered}{rgb}{.647,.129,.149}
\definecolor{mutedgrey}{rgb}{0.4,0.4,0.4}
\definecolor{shadecolor}{cmyk}{0,0,0.25,0.07}
\definecolor{MyDarkBlue}{rgb}{0,0.08,0.45}
\definecolor{MarginRed}{rgb}{0.8,0.0,0.0}
\definecolor{MarginBlue}{rgb}{0.2,0.0,1.0}
\definecolor{MarginGrey}{rgb}{0.4,0.4,0.4}
%\renewcommand{\chapnumfont}{\bfseries\Huge\sffamily}
%\renewcommand{\chaptitlefont}{\bfseries\Large\sffamily}
\setsecheadstyle{\bfseries\Large\sffamily\raggedright}
\setsubsecheadstyle{\bfseries\large\sffamily\raggedright}
\setsubsubsecheadstyle{\bfseries\normalsize\sffamily\raggedright}
\renewcommand \thesection{\bfseries\arabic{section}}
\makeatletter 
\newcommand\addRevisionData{%
\begin{picture}(0,0)% 
    \put(-110,-5){%
        \tiny% 
        {%
        {Published \today \enspace \copyright~Dale J. Brugh 
        }
}% 
}%
\end{picture}%
}
\flushbottom
\setstocksize{11in}{8.5in}
%\setlength{\parskip}{5pt}
\settrims{0pt}{0pt}
%\settrimmedsize{11in}{210mm}{*}
%\setlength{\trimtop}{0pt}
%\setlength{\trimedge}{\stockwidth}
%\addtolength{\trimedge}{-\paperwidth}
\settypeblocksize{8.5in}{5.0in}{*}
\setulmargins{1.25in}{*}{*}
\setlrmargins{1.25in}{*}{*}
\setmarginnotes{5mm}{4.0cm}{\onelineskip}
\setheadfoot{\onelineskip}{4\onelineskip}
\setheaderspaces{*}{\onelineskip}{*}
\checkandfixthelayout
%\setlength \fboxsep{0.1in}
\setlength \headwidth{\textwidth+\marginparwidth+\marginparsep}
%
\makepagestyle{courseinformation}
\makerunningwidth{courseinformation}{\headwidth}
\makeheadrule{courseinformation}{\headwidth}{\normalrulethickness}
\makeheadposition{courseinformation}{flushright}{flushleft}{flushleft}{flushleft}
\makeoddhead{courseinformation}%
    {\sffamily Course Syllabus: Chemistry 352 / Spring 2018}{}{\sffamily\thepage}

    \makeevenfoot{courseinformation}{}{}{}
    \makeoddfoot{courseinformation}{}{}{}
\makepagestyle{courseinformationtitle}
\makerunningwidth{courseinformationtitle}{\headwidth}
\makeheadposition{courseinformationtitle}{flushright}{flushleft}{flushleft}{flushleft}
    \makeevenfoot{courseinformationtitle}{}{}{}
    \makeoddfoot{courseinformationtitle}{}{}{}
\pagestyle{courseinformation}
\captionsetup{labelsep=colon,aboveskip=0.25cm,justification=RaggedRight,singlelinecheck=false,labelfont={bf,sf}}
%
% --- MARGIN FIGURE COMMAND ---
%
\newcommand{\marginfigures}[4]{
\marginpar{\centering
\includegraphics[width=#1]{#2}
\captionsetup{labelsep=newline,aboveskip=-0.5cm,justification=RaggedRight,singlelinecheck=false,labelfont={bf,sf}}
\captionsetup[figure]{position=bottom}
\captionof{figure}{#3}
\label{#4}}%
}%
%
% --- MARGIN NOTE COMMAND ---
%
\newcommand{\marginnote}[2]
{%
\marginpar{\raggedright\vspace{#1}\begin{Spacing}{0.65}\sffamily{{\tiny$\blacktriangleright$~\scriptsize#2}}\end{Spacing}} %
}
%
% --- SET UP TITLE ---
%
\setlength{\droptitle}{0.0in}
\backmatter
\pretitle{\noindent\huge\sffamily Course Syllabus \LARGE\par\noindent} 
\posttitle{\par\vskip 2.0em}
\preauthor{}
\postauthor{\par}
\predate{}
\postdate{\noindent\rule{\linewidth}{0.3pt}}
\title{Chemistry 352 / Spring 2018}
\date{}
\author{}
%
% --- BEGIN DOCUMENT ---
%
\begin{document}
\setsecnumdepth{subsubsection}
\maketitle
%\setsecnumdepth{subsection}
\thispagestyle{courseinformationtitle}
%
% --- INSTRUCTOR ---
%
\section{Instructor}
\begin{tabular}{rl|rl}
Name: & Dr. Dale J. Brugh & Office Phone: & 740-368-3530 \\
Email: & \href{mailto:djbrugh@owu.edu}{djbrugh@owu.edu} & Cell Phone: & 614-746-2397 \\
Office: & SCSC 262 & Office Hours: & Posted on Website \\
\end{tabular}
%
% --- MEETINGS ---
%
\section{Meetings}
\begin{tabular}{crcrl}
M & 1:10 p.m. & to & 4:00 p.m. & SCSC 355 \\
\end{tabular}

\section{Course Website}
The course website is a Blackboard course at \href{https://bb.owu.edu}{bb.owu.edu}. It is named Physical Chemistry Laboratory with course ID 2018SPRINGCHEM3521. The website is an important extension of the syllabus and should be read carefully.

\section{Laboratory Materials}
The following items will be used in this laboratory. 
\begin{enumerate}[1.]
\item Wolfram's \emph{Mathematica}, which is free for current OWU students from \href{https://user.wolfram.com/portal/login.html}{Wolfram's User Portal} at \href{https://user.wolfram.com/portal/login.html}{https://user.wolfram.com/portal/login.html}. 
\item \emph{Physical Chemistry} by Thomas Engel and Phillip Reid, Third Edition, Prentice Hall (2013).  ISBN-13: 978-0321812001
\item \emph{Quanta, Matter, and Change} by Peter Atkins, Julio de Paula, and Ronald Friedman, First Edition, W.H. Freeman and Company (2009). ISBN-10: 07167-6117-3, ISBN-13: 978-0-7167-6117-4.
\end{enumerate}

\section{Laboratory Prerequisites}
To take this laboratory course you must have passed one semester of physical chemistry that covers quantum mechanics. The easiest way to meet this prerequisite is to have credit for CHEM350. 

\section{Laboratory Content}
Chemistry 352 is an introduction to mathematical and molecular modeling methods in chemistry. The emphasis is on answering chemical questions using available tools. The topics to be covered are listed in Table~\ref{tab:topics}. These may not be covered in this order, and not all topics will be covered with equal depth. 

\begin{table}[h]
\caption{\sffamily Topics covered in Chemistry 352.}
\label{tab:topics}
\begin{tabular}{l|l} \toprule
Using Mathematica & Bonding \\
Linear Variational Theory & H\"uckel Theory \\
Atomic Property Calculation & Basis Sets  \\
Energy Calculations & Geometry Optimization \\
\bottomrule
\end{tabular}
\end{table}
%
% --- LABORATORY GOALS ---
%
\section{Laboratory Goals}
This laboratory has several goals. Each action I take is designed to move us closer to achieving one or all of these goals. My goals in this laboratory are
\begin{enumerate}[1.]
\item To provide you with an overview of molecular modeling and electronic structure techniques for computing properties of molecules,
\item To give you practice carrying out molecular modeling and electronic structure calculations on molecules using a research-grade tool,
\item To make you better able to recognize, understand, and interpret molecular modeling and electronic structure results presented in the literature,
\item To make you better able to formulate chemical questions that can be answered with molecular modeling and electronic structure methods, and
\item To provide you with the tools to continue your education in chemistry in the very best graduate programs on planet Earth.
\end{enumerate}
%
% --- LABORATORY EXPECTATIONS ---
%
\section{Laboratory Expectations}
You are expected to read the assigned readings and to come to each laboratory meeting prepared to participate. You are expected to display a high level of independence and a strong willingness to solve problems through your own troubleshooting. In short, you are expected to be professional. It is assumed that you are passionate about chemistry.
%
% --- WEEKLY ROUTINE ---
%
\section{Weekly Routine}
During \marginnote{-0.1in}{All projects are completed as part of an assigned team.} laboratory meetings you will work on projects intended to illustrate particular concepts in computational chemistry and to challenge your problem solving abilities. Projects will span one or two weeks. A pre-lab reading may be assigned for each project. At the end of each project you will submit a Mathematica notebook. A detailed schedule can be found on the course website.
%
% --- THINGS I GRADE ---
%
\section{Things I Grade}
You and I will determine your progress in this laboratory using the scores derived from evaluating your Mathematica notebooks. No other items are evaluated and graded. 
%
% --- Project Reports ---
%
\subsection{Project Notebooks}
For \marginnote{-0.1in}{Each team submits one Mathematica notebook, and all members of the team receive the same score.} each laboratory project, you must complete a Mathematica notebook containing your work. It will include justifications for the work, clear setup of the work, well-organized conclusions, and explanations of results. Notebooks will be evaluated for clarity, presentation, and accuracy. Each notebook is evaluated on a 100-point scale, and all 
%
% --- COURSE LETTER GRADE ---
%
\section{Course Letter Grade}
Your \marginnote{-0.1in}{Evaluation Weights \\ Project Notebooks: 100\% } course score will be calculated as a weighted average of the scores you earn on evaluation of the project notebooks ($100\%$). 
%
% --- LETTER GRADE ---
%
\section{Letter Grade}
Letter grades are assigned at the end of the course according to the minimum course score requirements listed in Table~\ref{tab:lettergrades}. Course scores below $55\%$ are considered failing. Please see \href{https://dephlogisticated.net/lettergrades}{dephlogisticated.net/lettergrades} for more detail about how your course letter grade is determined. 
\begin{table}[h]
\caption{\sffamily Minimum course scores necessary for each letter grade.}
\label{tab:lettergrades}
\begin{tabular}{cl||cl} \toprule
\textbf{Minimum Score} & \textbf{Letter Grade} & \textbf{Minimum Score} & \textbf{Letter Grade} \\ \hline
97 & \hspace{0.3in}A$+$ & 72 & \hspace{0.3in}C$+$ \\
88 & \hspace{0.3in}A & 68 & \hspace{0.3in}C \\
85 & \hspace{0.3in}A$-$ & 65 & \hspace{0.3in}C$-$ \\
82 & \hspace{0.3in}B$+$ & 62 & \hspace{0.3in}D$+$ \\
78 & \hspace{0.3in}B & 58 & \hspace{0.3in}D \\
75 & \hspace{0.3in}B$-$ & 55 & \hspace{0.3in}D$-$ \\
\bottomrule
\end{tabular}
\end{table}
%
% --- ADDITIONAL INFORMATION ---
%
\section{Additional Information}
Please see the course website for additional information. 
%
% ---LICENSE AND SOURCE CODE ---
%
\section{License and Source Code}
\copyright\ 2018 by Dale J. Brugh. Course Syllabus for Chemistry 352 (2018) is made available under a Creative Commons Attribution-ShareAlike 4.0 International License (CC BY-SA 4.0). See \href{https://creativecommons.org/licenses/by-sa/4.0/}{https://creativecommons.org/licenses/by-sa/4.0/} for details about your rights under this license. The \LaTeX\ source code is available under the same license from Github at \href{https://github.com/djbrugh/chem352-syllabus}{https://github.com/djbrugh/chem352-syllabus}. \marginpar{\vspace{-0.85in}\hfill\noindent\href{https://creativecommons.org/licenses/by-sa/4.0/}{\includegraphics[width=1in]{figs/cc-by-sa.pdf}}}


\end{document}